\documentclass[11pt,a4paper]{article}
\usepackage[margin={2.75cm,2.75cm}]{geometry}

\usepackage[spanish,es-tabla]{babel}
\usepackage[utf8]{inputenc}
\usepackage[T1]{fontenc}

\usepackage{hyperref}


\title{\LARGE \textbf{Práctica 1: Introducción a R y RStudio} \\[0.5em]
\Large Laboratorio de Matemáticas \\
\small \ \\
Doble Grado en Física y Matemáticas - Grado en Ingeniería Matemática}
\author{Miguel Beltrá Vidal}
\date{Curso 2026/27}


\begin{document}

\maketitle

\section{Objetivos}

El objetivo de esta práctica es el de introducir el software R al estudiante, así como el entorno de desarrollo RStudio. Después de esta práctica, el estudiante estará familiarizado con el software de forma que podrá ejecutar sus primeros comandos y escribir sus primeros programas en R.


\section{Desarrollo de la práctica}

\subsection{Un poco de historia...}

\subsubsection{¿Qué es R?}

De acuerdo con la página oficial (\url{https://www.r-project.org/}):

\begin{center}
\emph{R is a free software environment for statistical computing and graphics.}
\end{center}


\noindent Si buscamos el apartado \texttt{About R}, se nos dice lo siguiente:


\begin{center}
\begin{minipage}{0.9\linewidth}
	\emph{R is a language and environment for statistical computing and graphics. It is a GNU project which is similar to the S language and environment which was developed at Bell Laboratories (formerly AT\&T, now Lucent Technologies) by John Chambers and colleagues. R can be considered as a different implementation of S. There are some important differences, but much code written for S runs unaltered under R.}
	
	\medskip

	\emph{R provides a wide variety of statistical (linear and nonlinear modelling, classical statistical tests, time-series analysis, classification, clustering, …) and graphical techniques, and is highly extensible. The S language is often the vehicle of choice for research in statistical methodology, and R provides an Open Source route to participation in that activity.}
	
	\medskip

	\emph{One of R’s strengths is the ease with which well-designed publication-quality plots can be produced, including mathematical symbols and formulae where needed. Great care has been taken over the defaults for the minor design choices in graphics, but the user retains full control.}

	\medskip

	\emph{R is available as Free Software under the terms of the Free Software Foundation’s GNU General Public License in source code form. It compiles and runs on a wide variety of UNIX platforms and similar systems (including FreeBSD and Linux), Windows and MacOS.}
\end{minipage}
\end{center}

Así pues, podemos entender R como un lenguaje de programación, es decir una serie de reglas que nos permiten escribir un programa, y también como un entorno para realizar cálculos y gráficos estadísticos. Fue inicialmente desarrollado por Ross Ihaka and Robert Gentleman, en 1993. Por su larga trayectoria, se considera uno de los softwares más importantes para el tratamiento estadístico de los datos.


\subsection{Instalación de R y RStudio}


\subsection{Primeros pasos}

\subsection{Ejercicios}

\section{Entrega}

\end{document}